\documentclass[10pt,landscape,a4paper]{article}
\usepackage[ngerman]{babel}
\usepackage[T1]{fontenc}
\usepackage[dvipsnames]{xcolor}
%\usepackage[LY1,T1]{fontenc}
\usepackage{tikz}
\usetikzlibrary{shapes,positioning,arrows,fit,calc,graphs,graphs.standard}
\usepackage[nosf]{kpfonts}
\usepackage{multicol}
\usepackage{wrapfig}
\usepackage[top=2mm,bottom=3mm,left=2mm,right=2mm,headsep=1mm,includehead]{geometry}
\usepackage[framemethod=tikz]{mdframed}
\usepackage{microtype}
\usepackage{pdfpages}
\usepackage{csquotes}
\usepackage{booktabs}

\usepackage{amsmath}
\usepackage{bbm}
\newcommand{\N}{\mathbb{N}}
\newcommand{\Z}{\mathbb{Z}}
\newcommand{\R}{\mathbb{R}}
\newcommand{\C}{\mathbb{C}}
\newcommand{\K}{\mathbb{K}}
\newcommand{\e}{\varepsilon}

\DeclareMathOperator{\HP}{HP}
\DeclareMathOperator{\ggT}{ggT}
\DeclareMathOperator{\kgV}{kgV}

\newcommand\mat[1]{
	\begin{pmatrix}
		#1
	\end{pmatrix}
}

\usepackage{blindtext}

\usepackage{titlesec}
\input{../layout.tex}

\usepackage{hyperref}
\newcommand\nohyper[1]{\begin{NoHyper}#1\end{NoHyper}}

\usepackage{lastpage,fancyhdr}
\pagestyle{fancy}
\fancyhf{}
\fancyhead[L]{Rasmus Buurman}
\fancyhead[C]{\textsc{Mathematik für Informatik III}}
\fancyhead[R]{Seite \thepage{}/\nohyper{\pageref{LastPage}}}
\renewcommand\headrulewidth{0.5pt}

\begin{document}
\small
\begin{multicols*}{5}
	\section*{Metrische Räume}

\subsection*{Norm}
$V$ ist $\K$-VR. Norm $\|\cdot\|: V\to\R_{\ge 0}$:
\begin{enumerate}
	\item $\|x\| = 0 \iff x = 0$
	\item $\|\lambda x\| = |\lambda| \cdot \|x\|$
	\item $\|x+y\| \le \|x\| + \|y\|$
\end{enumerate}
Normen:
\begin{itemize}
	\item $\|x\|_2 := \sqrt{|x_1|^2 + \ldots + |x_d|^2}$
	\item $\|x\|_\infty := \max|x_i|$
\end{itemize}
Äquivalent, wenn: $c\|x\| \le \|x\|' \le C\|x\|$

\subsubsection*{Eigenschaften}
\begin{itemize}
	\item $\|x\| \ge 0$
	\item $\|x-y\| \ge |(\|x\| - \|y\|)|$
	\item Normen auf $\R^n$ sind äquivalent
\end{itemize}

\subsection*{Skalarprodukt}
$V$ ist $\K$-VR. S.produkt $\langle\cdot,\cdot\rangle : V\times V \to \K$:
\begin{enumerate}
	\item $\langle x, \lambda y + \mu z\rangle = \lambda \langle x,y\rangle + \mu \langle x, z\rangle$
	\item $\langle x,y\rangle = \overline{\langle y,x \rangle}$
	\item $\langle x,x \rangle\in\R$ und $\langle x,x \rangle > 0$ für $x \neq 0$
\end{enumerate}
$\|x\| := \sqrt{\langle x, x \rangle}$ ist Norm.

\subsubsection*{Cauchy-Schwarz-Ungleichung}
\[|\langle x,y \rangle| \le \|x\| \cdot \|y\|\]
Gleichheit, wenn $x$ und $y$ lin. unabh.

\subsection*{Metrik}
$M$ nicht-leer. Metrik $d: M \times M \to \R$:
\begin{enumerate}
	\item $d(x,y)=0 \iff x=y$
	\item $d(x,y) = d(y,x)$
	\item $d(x,y) \le d(x,z) + d(z,y)$
\end{enumerate}
\subsubsection*{Eigenschaften}
\begin{itemize}
	\item $d(x,y) \ge 0$
	\item $d(x,y) := \|x-y\|$ ist Metrik
\end{itemize}

	\section*{Topologische Grundbegriffe}

\subsection*{$\e$-Umgebung}
$M$ ist MR, $a\in M$, $\e > 0$.
\[U_\e(a) := \{x \in M \mid d(x,a) < \e\}\]
\subsubsection*{Eigenschaften}
\begin{itemize}
	\item $x \neq y \Rightarrow \exists \e: U_\e(x) \cap U_\e(y) = \emptyset$
	\item $\e$-Umgebungen sind offen
\end{itemize}

\subsection*{Offen \& Abgeschlossen}
$M$ ist MR. $U,O,A \subseteq M$ und $a \in U$.

Innerer Punkt $a$: $\exists \e>0: U_\e(a) \subseteq U$.

$O$ offen: Jeder Pkt in $O$ ist innerer Pkt.

$A$ abgeschlossen: $M \setminus A$ offen.

\subsection*{Abschluss \& Inneres}
$M$ ist MR, $U \subseteq M$.

Abschluss $\bar U$: Schnitt abg. Teilm. mit $U$

Inneres $\mathring U$: Vereinigung off. Teilm. mit $U$

\subsubsection*{Eigenschaften}
\begin{itemize}
	\item $\mathring U$ sind alle inneren Punkte
	\item $\bar U = U \cup \partial U = U \cup \HP(U)$
\end{itemize}

\subsection*{Häufungs- \& Randpunkt}
$M$ ist MR, $U \subseteq M, a \in M$.

\subsubsection*{Häufungspunkt $a \in \HP(U)$}
$\forall \e: \exists x \in U\setminus\{a\}: 0 < d(x,a) < \e$

bzw.: $\forall \e: U_\e(a) \cap (U \setminus \{a\}) \neq \emptyset$

\subsubsection*{Randpunkt $a \in \partial U$}
$\forall \e: U_\e(a) \cap U \neq \emptyset \land U_\e(a) \cap (M \setminus U) \neq \emptyset$

\subsection*{Beschränkt \& Kompakt}
$M$ ist MR, $U \subseteq M$.

\subsubsection*{$U$ ist beschränkt}
\[\exists r>0, a\in M: \forall x\in U: d(x,a) \le r\]

\subsubsection*{$(U_i)_{i\in I}$ ist offene Überdeckung}
$U_i$ offen und $U \subseteq \bigcup_{i\in I} U_i$

\subsubsection*{$U$ ist kompakt}
Jede offene Überdeckung von $U$ besitzt endliche Teilüberdeckung.

\subsubsection*{Eigenschaften}
\begin{itemize}
	\item Kompakt $\implies$ abgeschl. \& beschr.
	\item $A \subseteq \R^n$: Kmpkt. $\Leftrightarrow$ abg. \& besch.
\end{itemize}

\subsection*{Folgen in metrischen Räumen}
$M$ ist MR. Folge $(a_n)_{n\in\N}$ mit $a_n\in M$.

\subsubsection*{Konvergenz mit Grenzwert $a$}
\[\forall \e: \exists n_\e: \forall n \ge n_\e: d(a_n, a) < \e\]

$a_n \to a \iff d(a_n,a) \to 0$

$a_n \to a \land b_n \to b \implies d(a_n,b_n) \to d(a,b)$.

\subsubsection*{Cauchy-Folge}
\[\forall \e: \exists n_\e: \forall m>n \ge n_\e: d(a_m, a_n) < \e\]
Konvergent $\implies$ Cauchy-Folge

\subsubsection*{Beschränkte Folge}
$(a_n)_{n\in\N}$ beschr., wenn $\{a_n \mid n\in\N\}$ beschr.

Cauchy- \& konv. Folgen sind beschränkt.

\subsubsection*{Komponentenfolge}
$(a_k)_{k\in\N} \in \R^n$ und $c \in \R^n$.

$a_k \to c_k \iff a_{ki} \to c_i$

	\section*{Mehrdim. Funktionen}

\subsection*{Stetigkeit}
$M,M'$ sind MR. $f: M \to M'$ stetig:

$\forall \e \exists \delta_\e:\forall x \in M: $ \\ $d(x,a) < \delta_\e \implies d'(f(x), f(a)) < \e$

\subsubsection*{Grenzwertkriterium}
$a \in \HP(M)$. $f$ stetig $\displaystyle \Leftrightarrow \lim_{x\to a}f(x)=f(a)$.

\subsubsection*{Folgenkriterium}
$f$ stetig genau dann, wenn

$\displaystyle \forall (a_n)_{\N}: \lim_{n\to\infty}a_n=a \Rightarrow \lim_{n\to\infty}f(a_n)=f(a)$

\subsubsection*{Vektorwertige Funktionen}
$f:D\subseteq\R^n \to \R^m$ ist stetig, wenn alle Komponenten stetig sind.

\subsection*{Satz von Weierstraß}
$D \subseteq \R^n$ kompakt, $f:D\to\R^m$ stetig.
\begin{itemize}
	\item $f(D) \subseteq \R^m$ ist kompakt
	\item $f$ nimmt Min. und Max. an
\end{itemize}

	\section*{Differentialrechnung}
Sei $D\subseteq \R^n$ offen, $a\in D$.

\subsection*{Jacobi-Matrix}
\[
	f'(a) := \mat{
		\frac{\partial f_1}{\partial x_1}(a) & \ldots & \frac{\partial f_1}{\partial x_n}(a) \\
		\vdots & \ddots & \vdots \\
		\frac{\partial f_m}{\partial x_1}(a) & \ldots & \frac{\partial f_m}{\partial x_n}(a)
	}
\]

\subsubsection*{Gradient}
Für skalare Funktionen: $\nabla f(a) := f'(a)^T$

\subsection*{Totale Differenzierbarkeit}
$f: D \to \R^m$ in $a$ total diff.bar, wenn
\begin{enumerate}
	\item $f$ in $a$ partiell diff.bar und
	\item $f(x)=f(a)+f'(a)\cdot(x-a) + R(x)$ mit $\lim_{x\to a}\frac{\|R(x)\|}{\|x-a\|}=0$
\end{enumerate}

\subsubsection*{Zusammenhang partielle Diff.barkeit}
Partielle Ableitungen von $f$ existieren und sind stetig
$\implies f$ ist (total) diff.bar

\subsection*{Differenzierbarkeit/Stetigkeit}
$f$ (total) diff.bar $\implies f$ stetig

\subsection*{Richtungsableitung}
$f: D \to \R, a\in D, v\in\R^n$ mit $\|v\|=1$.

\[\frac{\partial f}{\partial v}(a) := \lim_{h\to0}\frac{f(a+h\cdot v) - f(a)}{h}\]

\subsubsection*{Zusammenhang Gradient}
\[\frac{\partial f}{\partial v}(a)=f'(a)\cdot v = \|f'(a)\| \cdot \cos\alpha\]

\subsection*{Mehrfaches partielles Ableiten}

\subsubsection*{Schwarz}
Ist $f:D\to\R$ $k$-mal stetig diff.bar, ist die Reihenfolge beim Ableiten egal.

\subsubsection*{Hesse-Matrix}
\[
	H_f(a) := \mat{
		\frac{\partial^2 f}{\partial x_1\partial x_1}(a) & \ldots & \frac{\partial^2 f}{\partial x_1 \partial x_n}(a) \\
		\vdots & \ddots & \vdots \\
		\frac{\partial^2 f}{\partial x_n\partial x_1}(a) & \ldots & \frac{\partial^2 f}{\partial x_n \partial x_n}(a)
	}
\]

\subsection*{Extremstellen}

\subsubsection*{Extremstellen sind kritische Punkte}
$a$ Extremstelle $\implies f'(a)=(0,\ldots,0)$

\subsubsection*{Hinreichende Bed. für Extremstellen}
$f$ 2-fach stetig diff.bar und $a$ kritisch.

Betrachte $H_f(a)$:
\begin{itemize}
	\item pos. definit $\implies$ iso. lok. Min.
	\item neg. definit $\implies$ iso. lok. Max.
	\item indefinit $\implies$ Sattelpunkt
\end{itemize}

	\section*{Satz von Taylor}

\subsection*{Eindimensional}
\[T_n(x) := \sum_{k=0}^n \frac{f^{(k)}(a)}{k!}(x-a)^k\]

\subsection*{Taylorpolynom}
\[T_{f,a}^k := \sum_{|\alpha| = 0}^{k}\frac{D^\alpha f(a)}{\alpha!} \cdot (t - a)^{\alpha}\]

\subsubsection*{Zweites Taylorpolynom $T_{f,a}^2(x)$}
\[f(a) + f'(a)(x-a) + \frac{1}{2}(x-a)^T H_f(a) (x-a)\]

	\section*{Umkehr- \& implizite Fnkt.}

\subsection*{Satz über die Umkehrfunktion}
$U\subseteq\R^n$ offen, $f:U\to\R^n$ stetig diff.bar.

$f'(x_0)$ inv.bar $\implies \exists U_0 \subseteq U$ offen mit:
\begin{itemize}
	\item $x_0\in U_0 \land \forall x\in U_0: \det(f'(x)) \neq 0$
	\item $f_{\mid_{U_0}}$ injektiv und $f(U_0)$ offen
	\item Umkehrfkt. $g$ von $f_{\mid_{U_0}}$ stg. diff.bar
\end{itemize}
$g'(f(x))=(f'(x))^{-1}$, $g'(x) = (f'(g(x)))^{-1}$

\subsection*{Implizite Funktionen}
$U \subseteq \R^{n+m}$ offen, $f: U \to \R^m$ stg. diff.bar.

Wenn $f(a,b)=0$ und $\det(f_y(a,b)) \neq 0$:

$\exists \varphi:U_\e(a) \to U_r(b)$ mit $f(x,\varphi(x))=0$

$\varphi'(x)=-(f_y(x,\varphi(x)))^{-1} \cdot f_x(x,\varphi(x))$

$\varphi'(a) = -(f_y(a,b))^{-1} \cdot f_x(a,b)$

\subsubsection*{Beispiel}
$f(a,b,c)=a^2+b^2+c^2-1=0$ \\
im Punkt $(a_0,b_0,c_0) = (0, \frac{1}{\sqrt{2}}, \frac{1}{\sqrt{2}})$

Suche Jacobi von $c=\varphi(a,b)$ in $(a_0,b_0)$.

Anwendbarkeit: \\
$f(a_0,b_0,c_0)=0$ und $f_c(a_0,b_0,c_0) \neq 0$.

$\varphi'(ab_0)=-(f_c(ab_0))^{-1} \cdot \big(f_a(ab_0) f_b(ab_0)\big)$ \\
$= - (\sqrt{2})^{-1} \cdot (0 ~ \sqrt{2}) = (0 ~ -1)$

	\section*{Extrema Nebenbedingungen}

\subsection*{Lagrange-Multiplikatoren}
\subsubsection*{Beispiel}
$f(x,y) = 3x+2y$ mit Nebenbedingung \\
$A = \{(x,y)^T \mid x^2 + y^2 = 1\}$

Eine Nebenbed.: $g(x,y)=1-x^2-y^2$ \\
(mehr sind aber auch möglich)

$F(x,y) = 3x + 2y + \lambda_1(1-x^2-y^2)$

Forderung: $F'(x,y) = 0$ \\
$\implies 3-2\lambda_1x = 0$ \\
$\implies 2-2\lambda_1y = 0$ \\
$\implies 1-x^2-y^2 = 0$

$a_1 = (\frac{3}{\sqrt{13}}, \frac{2}{\sqrt{13}})^T$,
$a_2 = (-\frac{3}{\sqrt{13}}, -\frac{2}{\sqrt{13}})^T$,

Gleichungssystem liefert mögliche Extremstellen.
Kriterium ist nötig, aber nicht hinreichend.

	\section*{Mehrdim. Integration}

\subsection*{Parameterintegrale}
$D=[a,b]\times[c,d]\subseteq\R^2$, $f:D\to\R$ stetig.

Dann existiert $F:[a,b]\to\R$ und stetig:
$F(x)=\int_c^d f(x,t)dt$

Wenn $\frac{\partial}{\partial x}f$ stetig:
$F'(x) = \int_c^d \frac{\partial}{\partial x}f(x,t)dt$

\subsubsection*{Satz von Fubini}
$\int_a^b F(x) = \int_a^b \int_c^d f(x,t)dt = \int_c^d \int_a^b f(x,t)dt$

\subsection*{Leibniz-Regel}
$f:[a,b]\times[c,d]\to\R$, $f$ und $f_x$ stetig.

$g,h: [a,b] \to [c,d]$ stetig diffbar.

Dann $F(x)=\int_{g(x)}^{h(x)}f(x,t)dt$ diffbar mit:

$F'(x) = \int_{g(x)}^{h(x)} f_x(x,t)dt + f(x,h(x))\cdot h'(x) - f(x,g(x))\cdot g'(x)$

\subsection*{Mehrfachintegrale}

Riemann-intbar: Ober- $=$ Unterintegral.

\subsubsection*{Iterierte Integrale}
Auf $I=[a,b]=[a_1,b_1]\times\ldots\times[a_n,b_n]\subseteq\R^n$ ist Integrationsreihenfolge egal.

	\section*{Euklidischer Algorithmus}

\subsection*{Beispiel: Euklidischer Algorithmus}
$a=48$ und $b=-30$
\begin{align*}
	48 &= (-1) \cdot (-30) + 18 \\
	-30 &= (-2) \cdot 18 + 6 \\
	18 &= 3 \cdot 6 + 0
\end{align*}
Als letzter Rest $\neq 0$ ist $\ggT(48,-30) = 6$.

\subsection*{Erweiterter euklidischer Algorithmus}

\subsubsection*{Lemma von B\'ezout}
$a,b\in\Z$, nicht beide 0.

Dann $\exists s,t\in\Z: \ggT(a,b)=sa + tb$.

\subsubsection*{Manuelle Methode}
\begin{enumerate}
	\item Berechne $\ggT(a,b)$
	\item Schrittw. rückw. einsetzen
\end{enumerate}

\subsection*{Anwendung}

\subsubsection*{Multiplikatives Inverse in $\Z_n$}
$0\neq z \in \Z_n$ hat Inv. $\iff \ggT(z,n)=1$

EEA liefert Inv. als $s$.

\subsubsection*{Lineare diophantische Gleichung}
\[c_0 = c_1\cdot x_1 + \ldots + c_n \cdot x_n\]
hat Lösung $\iff \ggT(c_1,\ldots,c_n) \mid c_0$.

	\section*{Primfaktorzerlegung}

\subsection*{Fundamentalsatz d. elem. Zahlentheorie}
Jedes $n \ge 2$ kann (bis auf Reihenfolge eindeutig) dargestellt werden als:
\[n = p_1^{e_1} \cdot \ldots \cdot p_k^{e_k},\]
wobei $p_i$ Primzahlen und $e_i\in\N$ sind.

\subsection*{Satz von Euklid}
Es gibt unendlich viele Primzahlen.

\subsection*{$\ggT$ und $\kgV$}
\begin{itemize}
	\item PFZ $\ggT(a,b)$ durch Min. der Exp.
	\item PFZ $\kgV(a,b)$ durch Max. der Exp.
\end{itemize}

	\section*{Chinesischer Restsatz}
$m_i\in\N$ pw. teilerfremd, $a_i\in\Z$.

Dann $\exists x, 0 \le x < M := \Pi m_i$ mit:

$x \equiv a_1 \pmod{m_1}, \ldots, x \equiv a_n \pmod{m_n}$

Lösung eindeutig modulo $M$.

\subsection*{Berechnung}
$M_i := \frac{M}{m_i}$. EEA für: $t_i m_i + s_i M_i = 1$.

Dann: $x = \big(\sum a_i s_i M_i\big) \mod M$.

\subsection*{Modulos zusammenfassen}
$m_1,m_2\in\N$ teilerfremd.

$x \equiv 1 \pmod{m_1} \land x \equiv 1 \pmod{m_2}$

$\implies x \equiv 1 \pmod{m_1\cdot m_2}$

	\section*{Eulersche $\varphi$-Funktion}

$\varphi(n) \equiv$ Anz. zu $n$ teilerfr. Zahlen in $[1,n]$

\subsection*{Berechnung}
$p$ Primzahl $\implies \varphi(p^e)=p^e-p^{e-1}$

$m_i$ pw. teilerfr. $\implies \varphi(\Pi m_i)=\Pi\varphi(m_i)$

\subsection*{Anwendung}

\subsubsection*{Kleiner Satz von Fermat}
$p$ Primzahl, $x$ teilerfremd zu $p$.
\[x^{p-1} \equiv 1 \pmod{p}\]

\subsubsection*{Eulerscher Satz}
$n\in\N$ und $a\in\Z$ teilerfremd.
\[a^{\varphi(n)} \equiv 1 \pmod{n}\]

	\section*{Polynomring}

\subsection*{Irreduzibilität von $p$}
$p = f \cdot g \implies \deg f=0 \lor \deg g = 0$

\subsubsection*{Faktorisierung}
$p \in K[x]$ lässt sich in eindeutige irreduzible normierte Polynome faktorisieren.

\subsection*{Restklassenring $K[x]_{m(x)}$}
Alle Polynome $p$ mit $\deg p < \deg m$.

	\section*{Zyklische Codes}

Generatorpoly. $g$, Kontrollpoly. $h$ mit
$g(x) \cdot h(x) = x^n + 1$

$C = \{f \cdot g \mid f \in \Z_2[x]_{x^n + 1}\} \subseteq \Z_2[x]_{x^n + 1}$

$\forall c \in C: h \cdot c = 0 \pmod{x^n+1}$

\subsection*{CRC}
$k := n - \deg(g)$

\subsubsection*{Codierung}
\begin{enumerate}
	\item Datenpoly. $M$ mit $\deg M < k$
	\item $M' = M \cdot x^{n-k}$
	\item $M' = Q \cdot g + R$
	\item $C = M' - R$
\end{enumerate}

\subsubsection*{Verifizierung}
Ist $C$ durch $g$ teilbar?


	\section*{RSA-Verschlüsselung}

\subsection*{Schlüsselerzeugung}
\begin{enumerate}
	\item Primzahlen $p,q$ und $n := p \cdot q$
	\item $\varphi(n) = (p-1)\cdot(q-1)$
	\item $e$ teilerfremd zu $\varphi(n)$
	\item $0<d<\varphi(n)$, $ed \equiv 1 \pmod{\varphi(n)}$
	\item Public: $(e,n)$, Private: $d$
\end{enumerate}

\subsection*{Verschlüsselung}
Nachricht $m$ mit $0 \le m < n$.

$c = m^e \mod n$

\subsection*{Entschlüsselung}
$m = c^d \mod n$

	\section*{Fourier Reihen}

\subsection*{Eulersche Formel}
$e^{i\varphi} = \cos\varphi + i \sin\varphi$

\subsection*{$T$-periodische Funktionen}
$f: \R \to \K$ mit $f(t + T) = f(t)$

\subsubsection*{Normierung: $2\pi$-periodisch}
$f(t)$ ist $T$-periodisch.

$\omega := \frac{2\pi}{T}$

$F(t) = f(\frac{t}{\omega})$ ist $2\pi$-periodisch.

\subsubsection*{Integration über $T$-periodische Funktion}
$\int_0^T f(t)dt = \int_a^{T+a}f(t)dt$

\subsection*{Trigonometrisches Polynom}

\subsubsection*{Sinus-Cosinus-Form}
$f(t) = \frac{a_0}{2} + \sum_{n=1}^N (a_n \cos(n\omega t) + b_n\sin(n\omega t))$

\subsubsection*{Exponentialform}
$f(t) = \sum_{k=-N}^{N} c_k e^{ik\omega t}$

\subsubsection*{Umrechnung}
$a_0=2c_0$, $a_n = c_n + c_{-n}$, $b_n=(c_n - c_{-n})i$

$c_0=\frac{a_0}{2}$, $c_k = \frac{a_k - ib_k}{2}$, $c_{-k}=\frac{a_k+ib_k}{2}$

\subsection*{Fourier-Reihe $S_f(t)$}
$S_f(t) &= \sum_{k=-\infty}^\infty c_k e^{ik\omega t}$

$= \frac{a_0}{2} + \sum_{n=1}^\infty (a_n\cos(n\omega t)+b_n\sin(n\omega t))$

\subsubsection*{Formeln von Euler-Fourier}
$a_n = \frac{2}{T} \int_0^T f(t)\cos(n\omega t)dt$

$b_n = \frac{2}{T} \int_0^T f(t)\sin(n\omega t)dt$

$c_k = \frac{1}{T} \int_0^T f(t) e^{-ik\omega t}dt$

\subsubsection*{Konvergenz}
Partialsummen $S_N(t)$ von $S_f(t)$

$\lim_{N\to\infty}\|f(t)-S_N(t)\|_2 = 0$

\end{multicols*}
\end{document}
